\documentclass[12pt]{article}
\usepackage[T1]{fontenc}
\usepackage[utf8]{inputenc}
\usepackage{lmodern}
\usepackage{textcomp}
\usepackage{enumitem}
\usepackage{geometry}
\geometry{margin=1in}
\begin{document}
\begin{center}
Answer Key - Science - K-12 Level Assessment 25 \
\small (Answers are ordered exactly as in the question paper.)
\end{center}

\section*{Multiple Choice}
\begin{enumerate}[label=\arabic*.]
\item \textbf{Answer:} B
\\ \textit{Options:} A) gas.; B) water.; C) wind.; D) clouds.
\\ \textit{Note:} Snow, rain, hail, and fog are all forms of water because they are different states or manifestations of H$_{2}$O in the atmosphere, occurring under various temperature and pressure conditions.
\item \textbf{Answer:} B
\\ \textit{Options:} A) The Moon is made of hot gases.; B) The Moon is covered with many craters.; C) The Moon has many bodies of liquid water.; D) The Moon has the ability to give off its own light.
\\ \textit{Note:} The Moon is covered with many craters due to its lack of atmosphere, which allows meteoroids and other celestial debris to collide with its surface without significant erosion or weathering.
\item \textbf{Answer:} B
\\ \textit{Options:} A) double pan balance; B) graduated cylinder; C) centimeter ruler; D) thermometer
\\ \textit{Note:} A graduated cylinder is used to measure the volume of liquids, making it suitable for accurately determining the amount of rainfall collected in a rain gauge.
\item \textbf{Answer:} D
\\ \textit{Options:} A) Use more than one control group.; B) Try a different experimental group.; C) Change the hypothesis if it is incorrect.; D) Repeat the procedure exactly the same each time.
\\ \textit{Note:} Repeating the procedure exactly the same each time minimizes variability and ensures that any differences in results can be attributed to the independent variable being tested, thereby increasing the reliability and validity of the experiment.
\item \textbf{Answer:} B
\\ \textit{Options:} A) reduce the energy requirements; B) adapt to exploit changed resources; C) interbreed with similar tropical organisms; D) emigrate before the change is complete
\\ \textit{Note:} The species' survival is most probable if it can adapt to exploit changed resources because adaptability to new environmental conditions and resource availability is crucial for sustaining populations in shifting ecosystems.
\end{enumerate}

\section*{True / False}
\begin{enumerate}[label=\arabic*.]
\setcounter{enumi}{5}
\item \textbf{Answer:} True
\\ \textit{Options:} A) True; B) False
\\ \textit{Note:} The statement is true because conduction is the process by which heat is transferred through direct contact, and the hot soup transfers heat to the metal spoon handle when they come into contact with each other.
\item \textbf{Answer:} True
\\ \textit{Options:} A) True; B) False
\\ \textit{Note:} Solar power is considered a renewable energy source because it harnesses energy from the sun, which is abundant and inexhaustible, and it has minimal negative environmental impact compared to fossil fuels, as it produces no greenhouse gas emissions during operation.
\item \textbf{Answer:} True
\\ \textit{Options:} A) True; B) False
\\ \textit{Note:} Burning coal releases pollutants such as sulfur dioxide and particulate matter, which can lead to long-term degradation of air and water quality, ultimately causing slow changes in ecosystems.
\item \textbf{Answer:} False
\\ \textit{Options:} A) True; B) False
\\ \textit{Note:} Leaves primarily function to carry out photosynthesis, which provides energy for the plant, rather than producing flowers and seeds.
\item \textbf{Answer:} False
\\ \textit{Options:} A) True; B) False
\\ \textit{Note:} Ocean waves do not move more easily than continental rock because the solid structure of continental rock is much more rigid and stable compared to the fluid motion of ocean waves, which primarily affect weather patterns through energy transfer rather than direct movement of land.
\end{enumerate}

\section*{Long Form Response}
\begin{enumerate}[label=\arabic*.]
\setcounter{enumi}{10}
\item \textbf{Answer:} The Sun appears larger than other stars primarily due to its proximity to Earth. While many stars are much larger than the Sun, they are located at vast distances, making them appear smaller and dimmer in the night sky. The brightness of the Sun also contributes to its perceived size; it is the brightest object in our sky, which enhances our perception of its size compared to distant stars. Additionally, human perception plays a role; we are used to comparing objects in our environment based on their brightness and size, leading us to perceive the Sun as the largest star, even though it is not the biggest in the universe.
\\ \textit{Note:} The Sun appears larger than other stars because its close proximity to Earth makes it appear brighter and more dominant in the sky, while distant stars, despite being larger, appear smaller and dimmer due to their vast distances from us.
\item \textbf{Answer:} The body's mechanisms for heat release during and after exercise are crucial for temperature regulation and overall performance. When we exercise, our muscles generate heat, raising our body temperature. To cool down, the body sweats, and as sweat evaporates from the skin, it removes heat, helping to maintain a stable internal temperature. Additionally, increased blood flow to the skin allows more heat to escape. By effectively regulating temperature, the body can enhance endurance and prevent overheating, which contributes to better athletic performance.
\\ \textit{Note:} correct because it explains how sweating and increased blood flow to the skin are essential mechanisms for dissipating heat generated during exercise, thereby maintaining internal temperature and enhancing performance.
\end{enumerate}

\vfill
\noindent\textit{End of Answer Key}
\end{document}